% GENERATED FILE. Edit canonical lessons, then run npm run generate:books.
% canonical-source-hash: fnv1a64:74bd9fb281ecd864
% canonical-lessons: MR-C118-chalis, MR-C118-pannas, MR-C118-sath, MR-C118-sattar, MR-C118-navvad

\chapter{Forty, Fifty, Sixty, Seventy, Ninety}
\label{ch:mr-forty-fifty-sixty-seventy-ninety}
\hlchaptermodality{118}

\begin{chapteropening}
\textbf{By the end of this chapter:} \emph{I can say forty, fifty, sixty, seventy and ninety.}

Complete the last lesson of chapter 118: I can say forty, fifty, sixty, seventy and ninety.
\end{chapteropening}

\section[\texorpdfstring{\mr{चाळीस} (chāḷīs)}{चाळीस (chāḷīs)}]{\mr{चाळीस} (chāḷīs) — forty}

\label{lesson:MR-C118-chalis}



\begin{quote}
Before the new one: say the Marathi for a sandal, then the Marathi for thirty.
\end{quote}

\subsection*{You'll want to know: \mr{चाळीस}}
\textbf{\mr{चाळीस}} — \emph{chāḷīs} — \textquotedblleft{}forty\textquotedblright{}.

\begin{cousinweb}
\textbf{\mr{चाळीस}} (\emph{chāḷīs}) is a tadbhava, a word inherited from Sanskrit and reshaped by sound change. It comes from \emph{catvāriṁśat}, \textquotedblleft{}forty\textquotedblright{} (CDIAL 4656). CDIAL notes that its \emph{ch-}, where Marathi would say \emph{ts-}, was taken from Hindi.
\end{cousinweb}

\begin{grammarlens}[title={What you've built}]
One more word for this chapter.
\end{grammarlens}

\subsection*{Your turn}
\begin{itemize}
\raggedright
  \item \emph{Say it:} \emph{chāḷīs}
  \item \emph{Say it:} \emph{chāḷīs}, once more
  \item \emph{Say it:} say \emph{chāḷīs}
  \item \emph{Recall:} say the Marathi for a sandal, then the Marathi for thirty, then say \emph{chāḷīs} again
\end{itemize}

\begin{tcolorbox}[breakable,colback=teal!4,colframe=teal!35!black,title={Before you move on}]
What does \mr{चाळीस} mean? (\textquotedblleft{}Forty\textquotedblright{}.) Say it once more.
\end{tcolorbox}
\section[\texorpdfstring{\mr{पन्नास} (pannās)}{पन्नास (pannās)}]{\mr{पन्नास} (pannās) — fifty}

\label{lesson:MR-C118-pannas}



\begin{quote}
Before the new one: say the Marathi for thirty, then the Marathi for forty.
\end{quote}

\subsection*{You'll want to know: \mr{पन्नास}}
\textbf{\mr{पन्नास}} — \emph{pannās} — \textquotedblleft{}fifty\textquotedblright{}.

\begin{cousinweb}
\textbf{\mr{पन्नास}} (\emph{pannās}) is a tadbhava, a word inherited from Sanskrit and reshaped by sound change. It comes from \emph{pañcāśat}, \textquotedblleft{}fifty\textquotedblright{} (CDIAL 7682).
\end{cousinweb}

\begin{grammarlens}[title={What you've built}]
One more word for this chapter.
\end{grammarlens}

\subsection*{Your turn}
\begin{itemize}
\raggedright
  \item \emph{Say it:} \emph{pannās}
  \item \emph{Say it:} \emph{pannās}, once more
  \item \emph{Say it:} say \emph{pannās}
  \item \emph{Recall:} say the Marathi for thirty, then the Marathi for forty, then say \emph{pannās} again
\end{itemize}

\begin{tcolorbox}[breakable,colback=teal!4,colframe=teal!35!black,title={Before you move on}]
What does \mr{पन्नास} mean? (\textquotedblleft{}Fifty\textquotedblright{}.) Say it once more.
\end{tcolorbox}
\section[\texorpdfstring{\mr{साठ} (sāṭh)}{साठ (sāṭh)}]{\mr{साठ} (sāṭh) — sixty}

\label{lesson:MR-C118-sath}



\begin{quote}
Before the new one: say the Marathi for forty, then the Marathi for fifty.
\end{quote}

\subsection*{You'll want to know: \mr{साठ}}
\textbf{\mr{साठ}} — \emph{sāṭh} — \textquotedblleft{}sixty\textquotedblright{}.

\begin{cousinweb}
\textbf{\mr{साठ}} (\emph{sāṭh}) is a tadbhava, a word inherited from Sanskrit and reshaped by sound change. It comes from \emph{ṣaṣṭi}, \textquotedblleft{}sixty\textquotedblright{} (CDIAL 12804).
\end{cousinweb}

\begin{grammarlens}[title={What you've built}]
One more word for this chapter.
\end{grammarlens}

\subsection*{Your turn}
\begin{itemize}
\raggedright
  \item \emph{Say it:} \emph{sāṭh}
  \item \emph{Say it:} \emph{sāṭh}, once more
  \item \emph{Say it:} say \emph{sāṭh}
  \item \emph{Recall:} say the Marathi for forty, then the Marathi for fifty, then say \emph{sāṭh} again
\end{itemize}

\begin{tcolorbox}[breakable,colback=teal!4,colframe=teal!35!black,title={Before you move on}]
What does \mr{साठ} mean? (\textquotedblleft{}Sixty\textquotedblright{}.) Say it once more.
\end{tcolorbox}
\section[\texorpdfstring{\mr{सत्तर} (sattar)}{सत्तर (sattar)}]{\mr{सत्तर} (sattar) — seventy}

\label{lesson:MR-C118-sattar}



\begin{quote}
Before the new one: say the Marathi for fifty, then the Marathi for sixty.
\end{quote}

\subsection*{You'll want to know: \mr{सत्तर}}
\textbf{\mr{सत्तर}} — \emph{sattar} — \textquotedblleft{}seventy\textquotedblright{}.

\begin{cousinweb}
\textbf{\mr{सत्तर}} (\emph{sattar}) is a tadbhava, a word inherited from Sanskrit and reshaped by sound change. It comes from \emph{saptati}, \textquotedblleft{}seventy\textquotedblright{} (CDIAL 13143).
\end{cousinweb}

\begin{grammarlens}[title={What you've built}]
One more word for this chapter.
\end{grammarlens}

\subsection*{Your turn}
\begin{itemize}
\raggedright
  \item \emph{Say it:} \emph{sattar}
  \item \emph{Say it:} \emph{sattar}, once more
  \item \emph{Say it:} say \emph{sattar}
  \item \emph{Recall:} say the Marathi for fifty, then the Marathi for sixty, then say \emph{sattar} again
\end{itemize}

\begin{tcolorbox}[breakable,colback=teal!4,colframe=teal!35!black,title={Before you move on}]
What does \mr{सत्तर} mean? (\textquotedblleft{}Seventy\textquotedblright{}.) Say it once more.
\end{tcolorbox}
\section[\texorpdfstring{\mr{नव्वद} (navvad)}{नव्वद (navvad)}]{\mr{नव्वद} (navvad) — ninety}

\label{lesson:MR-C118-navvad}



\begin{quote}
Before the new one: say the Marathi for sixty, then the Marathi for seventy.
\end{quote}

\subsection*{You'll want to know: \mr{नव्वद}}
\textbf{\mr{नव्वद}} — \emph{navvad} — \textquotedblleft{}ninety\textquotedblright{}.

\begin{cousinweb}
CDIAL lists \textbf{\mr{नव्वद}} (\emph{navvad}) under Sanskrit \emph{navati}, \textquotedblleft{}ninety\textquotedblright{}, and marks the Marathi form as a borrowing (CDIAL 6995).
\end{cousinweb}

\begin{grammarlens}[title={What you've built}]
One more word for this chapter.
\end{grammarlens}

\subsection*{Your turn}
\begin{itemize}
\raggedright
  \item \emph{Say it:} \emph{navvad}
  \item \emph{Say it:} \emph{navvad}, once more
  \item \emph{Say it:} say \emph{navvad}
  \item \emph{Recall:} say the Marathi for sixty, then the Marathi for seventy, then say \emph{navvad} again
\end{itemize}

\begin{tcolorbox}[breakable,colback=teal!4,colframe=teal!35!black,title={Before you move on}]
What does \mr{नव्वद} mean? (\textquotedblleft{}Ninety\textquotedblright{}.) Say it once more.
\end{tcolorbox}
